% Part: incompleteness
% Chapter: arithmetization-syntax
% Section: substitution

\documentclass[../../../include/open-logic-section]{subfiles}

\begin{document}

\olfileid{inc}{art}{sub}
\olsection{جانشانی}

به یاد آورید جانشانی عملِ جایگزین‌کردن همهٔ رخدادهای آزادِ
!!a{variable}~$u$ در !!a{formula}~$!A$ با ترم~$t$ است و با
$\Subst{!A}{t}{u}$ نوشته می‌شود. وقتی این عمل روی اعداد گودلِ
!!{variable}s، !!{formula}s و ترم‌ها انجام شود، بازگشتی اولیه است.

\begin{prop}\ollabel{prop:subst-primrec}
تابع بازگشتی اولیه‌ای چون~$\fn{Subst}(x,y,z)$ وجود دارد با این ویژگی
که
\[
\fn{Subst}(\Gn{!A}, \Gn{t}, \Gn{u}) = \Gn{\Subst{!A}{t}{u}}.
\]
\end{prop}

\begin{proof}
می‌توانیم تابع~$\fn{hSubst}$ را با بازگشت اولیه چنین تعریف کنیم:
\begin{multline*}
\begin{aligned}
\fn{hSubst}(x, y, z, 0) & = \emptyseq \\
\fn{hSubst}(x, y, z, i+1) & =
\end{aligned}\\
\begin{cases}
\fn{hSubst}(x, y, z, i) \concat y & \text{اگر $\fn{FreeOcc}(x, z, i)$} \\
\fn{append}(\fn{hSubst}(x, y, z, i), (x)_{i}) & \text{در غیر این صورت.}
\end{cases}
\end{multline*}
اکنون می‌توان $\fn{Subst}(x,y,z)$ را برابر
$\fn{hSubst}(x,y,z,\len{x})$ تعریف کرد.
\end{proof}

\begin{prop}
\ollabel{prop:free-for}
رابطهٔ~$\fn{FreeFor}(x,y,z)$، که برقرار است اگر و تنها اگر ترم دارای
عدد گودل~$y$ برای متغیر دارای عدد گودل~$z$ در فرمول دارای عدد
گودل~$x$ !!{free for} باشد، بازگشتی اولیه است.
\end{prop}

\begin{proof} تمرین. \end{proof}

\begin{prob}
\olref[inc][art][sub]{prop:free-for} را ثابت کنید.
\end{prob}

\end{document}
